\documentclass[11pt]{article}
\usepackage[a4paper,margin=1in]{geometry}
\usepackage{amsmath,amssymb,amsthm,mathtools}
\usepackage{graphicx}
\usepackage[hidelinks]{hyperref}
\usepackage{microtype}
\usepackage[T1]{fontenc}
\usepackage{lmodern}
\setlength{\emergencystretch}{3em}

\newtheorem{theorem}{Theorem}
\newtheorem{lemma}{Lemma}
\newtheorem{corollary}{Corollary}
\newtheorem{remark}{Remark}
\newcommand{\N}{\mathbb N}
\newcommand{\R}{\mathbb R}
\newcommand{\eps}{\varepsilon}
\newcommand{\Rrho}{{}^{\rho}\widetilde{\mathbb R}}
\newcommand{\Ru}{{}^{\rho}\widetilde{\R}_{\mathrm u}}
\newcommand{\Rs}{{}^{\rho}_{\sigma}\widetilde{\R}_{\mathrm s}}
\newcommand{\Ns}{{}^{\sigma}\widetilde{\N}}

\title{The Natural Map to the Hyperseries Sequence Space Is Not Injective\\
\large A supermoderate-spike counterexample to the open problem in arXiv:2006.16141}
\author{Candidate full solution packet for arXiv:2006.16141}
\date{27 August 2026}

\begin{document}
\maketitle

\noindent\textbf{Status.}
Candidate full negative answer, likely valid and awaiting human review.  For
every pair of gauges for which the source map is defined, its kernel is
nonzero; in fact, the kernel contains infinitely many nonzero pairwise
orthogonal idempotents.

\section{Source question}

The source is Diksha Tiwari and Paolo Giordano, \emph{Hyperseries in the
non-Archimedean ring of Colombeau generalized numbers}, arXiv:2006.16141,
published in \emph{Monatshefte f\"ur Mathematik} 197 (2022), 193--223
\cite{TG}.  On source PDF page 7, immediately after Lemma 7(iii), the authors
construct an $\Rrho$-linear map
\[
 \lambda:\Ru\longrightarrow\Rs,
 \qquad [a_{n\eps}]_{\mathrm u}\longmapsto[a_{n\eps}]_{\mathrm s},
\]
under the gauge condition $\sigma\geq\rho^*$, and state:
\begin{quote}
``It remains an open problem the study of injectivity of the map $\lambda$.''
\end{quote}

\begin{figure}[ht]
  \centering
  \includegraphics[width=\textwidth]{figures/open_problem_crop.png}
  \caption{The injectivity question in the source paper, source PDF page 7.}
\end{figure}

\section{Definitions used}

Let $\rho=(\rho_\eps)$ and $\sigma=(\sigma_\eps)$ be gauges, so in
particular $\rho_\eps,\sigma_\eps>0$ and both tend to zero.  The condition
$\sigma\geq\rho^*$ means that
\[
 \sigma_\eps\geq\rho_\eps^Q
 \quad\hbox{eventually for some }Q\in\N_{>0}.
\]
An array $(a_{n\eps})_{n\in\N,\eps}$ is \emph{uniformly moderate} if
\[
 \exists Q\in\N\ \forall^{0}\eps\ \forall n\in\N:
 |a_{n\eps}|\leq\rho_\eps^{-Q},
\]
and is \emph{uniformly negligible} if
\[
 \forall q\in\N\ \forall^{0}\eps\ \forall n\in\N:
 |a_{n\eps}|\leq\rho_\eps^q.
\]
The quotient is $\Ru$.

Every $N=[N_\eps]\in\Ns$ has a $\sigma$-moderate representative: for
some $R\in\N$,
\[
                         N_\eps\leq\sigma_\eps^{-R}
                         \quad\hbox{eventually}.
\]
The target $\Rs$ identifies two arrays if every interval sum between
$\sigma$-moderate hypernatural endpoints is $\rho$-negligible.  Thus an
array represents zero in $\Rs$ precisely when, for all $M,N\in\Ns$,
\[
 \left(\sum_{n=N_\eps}^{M_\eps}a_{n\eps}\right)_\eps\sim_\rho0.
\]
These are the definitions in Section 2 of \cite{TG} specialized to the
zero class.

\section{Proof intuition}

The two quotient relations observe different ranges of the ordinary index
$n$.  Uniform negligibility in $\Ru$ tests \emph{every} $n\in\N$ at each
small $\eps$.  In contrast, $\Rs$ only tests intervals whose endpoints grow
at most like a fixed power of $1/\sigma_\eps$.

Place a single unit spike at an index $K_\eps$ growing faster than every such
power.  The domain still sees the spike, because its definition takes a
supremum over all ordinary $n$.  Every admissible hyperseries interval misses
it eventually, so the target sees the zero array.  The counterexample is the
gap between ``all ordinary indices'' and ``all $\sigma$-moderate indices''
made literal.

\section{Formal result}

\begin{lemma}[A super-$\sigma$ index]\label{lem:super}
Define
\[
                   K_\eps:=\left\lceil\exp(1/\sigma_\eps)\right\rceil.
\]
Then $K_\eps>\sigma_\eps^{-R}$ eventually for every fixed $R\in\N$.
Consequently, if $N=[N_\eps]\in\Ns$, then $N_\eps<K_\eps$ eventually.
\end{lemma}

\begin{proof}
Put $x_\eps=1/\sigma_\eps$.  Then $x_\eps\to+\infty$, and for every
fixed $R$,
\[
 \frac{\exp(1/\sigma_\eps)}{\sigma_\eps^{-R}}
   =\frac{e^{x_\eps}}{x_\eps^R}\longrightarrow+\infty.
\]
This proves the first assertion.  For the second, choose $R$ such that
$N_\eps\leq\sigma_\eps^{-R}$ eventually and apply the first assertion.
\end{proof}

\begin{theorem}[Complete negative answer to the injectivity problem]
\label{thm:main}
For every pair of gauges $\rho,\sigma$ satisfying $\sigma\geq\rho^*$, the
natural linear map
\[
                        \lambda:\Ru\longrightarrow\Rs
\]
of Lemma 7(iii) in \cite{TG} is not injective.
\end{theorem}

\begin{proof}
Let $K_\eps$ be as in Lemma \ref{lem:super} and define
\[
 a_{n\eps}:=
 \begin{cases}
  1,&n=K_\eps,\\
  0,&n\neq K_\eps.
 \end{cases}
\]
Because $|a_{n\eps}|\leq1=\rho_\eps^0$, this array is uniformly moderate.
It is not uniformly negligible: since $\rho_\eps\to0$, eventually
$\rho_\eps<1$, while for the ordinary index $n=K_\eps$ one has
$|a_{K_\eps,\eps}|=1>\rho_\eps$.  Therefore
\[
                         [a_{n\eps}]_{\mathrm u}\neq0
                         \quad\hbox{in }\Ru.
\]

Now fix arbitrary $M=[M_\eps]$ and $N=[N_\eps]$ in $\Ns$.  By Lemma
\ref{lem:super}, both endpoints are eventually strictly smaller than
$K_\eps$.  Hence every ordinary integer between the two endpoints is also
smaller than $K_\eps$, and
\[
                         \sum_{n=N_\eps}^{M_\eps}a_{n\eps}=0
                         \quad\hbox{eventually}.
\]
The interval-sum net is therefore $\rho$-negligible.  Since $M,N$ were
arbitrary, $(a_{n\eps})$ represents zero in $\Rs$.  Thus
\[
       \lambda([a_{n\eps}]_{\mathrm u})=[a_{n\eps}]_{\mathrm s}=0
       \quad\hbox{while}\quad [a_{n\eps}]_{\mathrm u}\neq0,
\]
which proves noninjectivity.
\end{proof}

\begin{corollary}[A large invisible tail]
The kernel of $\lambda$ contains infinitely many nonzero pairwise orthogonal
idempotents of the pointwise ring $\Ru$.
\end{corollary}

\begin{proof}
For $j\in\N_{>0}$ set
\[
 K^{(j)}_\eps=\left\lceil e^{j/\sigma_\eps}\right\rceil,
 \qquad
 a^{(j)}_{n\eps}=\mathbf 1_{\{K^{(j)}_\eps\}}(n).
\]
The proof of Theorem \ref{thm:main} applies to each $a^{(j)}$.  Each domain
class is nonzero and idempotent.  If $j\neq k$, then
$K^{(j)}_\eps\neq K^{(k)}_\eps$ eventually, so the pointwise product of the
two arrays is eventually identically zero, uniformly in $n$.  Their classes
are therefore orthogonal.
\end{proof}

\section{Verification notes}

\begin{itemize}
\item \textbf{Gauge generality.}  No special form, monotonicity, or comparison
beyond the source hypothesis $\sigma\geq\rho^*$ is used.  In fact, that
comparison is needed only so that the source map $\lambda$ is defined.
\item \textbf{Domain nonzero.}  The witness to failure of uniform
negligibility is allowed to depend on $\eps$: the quantifier order is
$\forall^0\eps\,\forall n$.  At each $\eps$, choosing the ordinary integer
$n=K_\eps$ contradicts the $q=1$ estimate.
\item \textbf{Target zero.}  Every individual hypernatural has some fixed
moderateness exponent.  After taking the maximum exponent for two endpoints,
Lemma \ref{lem:super} puts the spike beyond the entire interval.
\item \textbf{Representatives.}  The proof checks the target equivalence for
arbitrary hypernatural classes by choosing arbitrary moderate
representatives.  Changing representatives does not affect the conclusion,
since the source definition already makes interval sums well defined.
\item \textbf{No computation.}  The argument is exact and uses only
$e^x/x^R\to\infty$; no numerical or symbolic dependency remains.
\end{itemize}

\section{Scope and novelty check}

The theorem answers the printed injectivity problem completely and
negatively for every gauge pair in the stated domain of the map.  It does not
alter the usefulness of $\lambda$ for transferring convergence statements;
it identifies an ``invisible tail'' that all $\sigma$-moderate hyperseries
intervals fail to detect.

A bounded novelty check was performed on 27 August 2026.  The four local run
indexes were searched for arXiv:2006.16141 and the phrases ``hyperseries'',
``injectivity of the map lambda'', ``uniformly moderate'', and
``supermoderate''.  Exact-phrase web searches found the published source,
which still states the problem open, and a 2025 University of Vienna thesis
that reproduces the same open statement.  The local arXiv sources of the
follow-up \emph{Hyper-power series and generalized real analytic functions}
(arXiv:2212.04757) \cite{TMG}, as well as citing arXiv sources 2111.15408,
2501.17364, and 2604.17904, were searched for ``injectiv'', ``lambda'', and
the two quotient-space notations; no resolution was found.  No paper or
packet containing this supermoderate-spike counterexample was located.  The
search is bounded rather than exhaustive, so novelty confidence is moderate
to high.

\section{Human-review recommendation}

Promote as a full negative answer if a reviewer confirms the quantifier order
in uniform negligibility and the standard definition that every element of
$\Ns$ has a representative bounded by one fixed power of
$1/\sigma_\eps$.  Those are stated explicitly in the source.  The remaining
argument is elementary, works for arbitrary admissible gauges, and has no
unproved dependency.

\begin{thebibliography}{9}
\bibitem{TG}
D.~Tiwari and P.~Giordano,
\emph{Hyperseries in the non-Archimedean ring of Colombeau generalized
numbers}, Monatshefte f\"ur Mathematik \textbf{197} (2022), 193--223;
arXiv:2006.16141,
\href{https://doi.org/10.1007/s00605-021-01647-0}{doi:10.1007/s00605-021-01647-0}.

\bibitem{TMG}
D.~Tiwari, A.~Mukhammadiev, and P.~Giordano,
\emph{Hyper-power series and generalized real analytic functions},
arXiv:2212.04757 (2022).
\end{thebibliography}

\end{document}
